\documentclass[sigconf]{acmart}

\usepackage{amsmath}
\usepackage{mathtools}
\usepackage{acro}
\usepackage{cleveref}
\usepackage{listings}
\usepackage{subfiles}
\usepackage{xr-hyper}

\DeclareAcronym{KG}{
  short = KG,
  long  = Knowledge Graph,
}
\DeclareAcronym{BKG}{
  short = BKG,
  long  = Bag Semantic Knowledge Graph,
  long-plural-form = Bag Semantic Knowledge Graphs,
}
\DeclareAcronym{RDF}{
  short = RDF,
  long  = Resource Description Framework,
}
\DeclareAcronym{BGP}{
  short = BGP,
  long  = Basic Graph Pattern,
}
\DeclareAcronym{CQ}{
  short = CQ,
  long  = Conjunctive Query,
  long-plural-form = Conjunctive Queries,
}
\DeclareAcronym{UCQ}{
  short = UCQ,
  long  = Unions of Conjunctive Queries,
  long-plural-form = Unions of Conjunctive Queries,
}
\DeclareAcronym{UCFQ}{
  short = UCFQ,
  long  = Unions of Conjunctive Federated Queries,
  long-plural-form = Unions of Conjunctive Federated Queries,
}
\DeclareAcronym{UGP}{
  short = UGP,
  long  = Union Graph Pattern,
}
\DeclareAcronym{IRI}{
  short = IRI,
  long  = Internationalized Resource Identifier,
}
\DeclareAcronym{ESA}{
  short = ESA,
  long  = Exhaustive Source Assignments,
}
\DeclareAcronym{GGP}{
  short = GGP,
  long  = Group Graph Pattern,
}
\DeclareAcronym{GP}{
  short = GP,
  long  = Graph Pattern,
}
\DeclareAcronym{BFR}{
  short = BFR,
  long  = Bag Federation Reduction,
}


\newtheorem{proposition}{Proposition}[section]
\newtheorem{definition}{Definition}[section]
\newtheorem{corollary}{Corollary}[section]
\newtheorem{lemma}{Lemma}[section]

\newcommand{\eval}[2]{\llbracket #1 \rrbracket^{#2}}
\DeclareMathOperator{\fmj}{mj}
\DeclareMathOperator{\fmu}{mu}
\DeclareMathOperator{\freq}{req}

\setcopyright{cc}
\copyrightyear{2026}
\acmYear{2026}
\acmDOI{10.1145/0000000000}
\acmConference[EDBT '26]{Extending Database Technology}{TBD}{TBD}
\acmISBN{978-1-4503-XXXX-X/26/03}

\title{Bag Semantics and Query Containment in SPARQL Federation}

\author{Bryan-Elliott Tam}
\orcid{0000-0003-3467-9755}
\email{bryan-elliott.tam@ugent.be}
\affiliation{%
  \institution{Ghent University}
  \city{Ghent}
  \country{Belgium}
}

\author{Pieter Colpaert}
\orcid{0000-0001-6917-2167}
\email{pieter.colpaert@ugent.be}
\affiliation{%
  \institution{Ghent University}
  \city{Ghent}
  \country{Belgium}
}

\author{Ruben Taelman}
\orcid{0000-0001-5118-256X}
\email{ruben.taelman@ugent.be}
\affiliation{%
  \institution{Ghent University}
  \city{Ghent}
  \country{Belgium}
}

\keywords{Linked Data, SPARQL, Query containment, Federated Querying,
  bag-semantic, bag-set semantic, multiset}

\begin{document}

\begin{abstract}
Querying over federations of distributed data sources enables applications to integrate heterogeneous datasets without centralizing data, but can impose significant computational overhead on both federation endpoints and querying clients.
In contexts such as scientific reproducibility studies, it is often necessary to re-execute queries or run refinement queries.
However, reissuing these queries can lead to unnecessary load, creating a barrier to such use cases.
To mitigate this, it is necessary to develop mechanisms that allow reusing previously computed query results to reduce redundant computation and network traffic.
In this paper, we explore query containment as a mechanism to enable optimization of federated queries.
The current literature on SPARQL query containment focuses on set semantics, whereas the SPARQL query language operates under bag-set semantics.
Furthermore, the query containment problem has not been explored in the context of federated queries.
This paper bridges this gap by exploring SPARQL query containment for real-world applications, covering both bag semantics and bag-set semantics for fragments of single endpoint and federated queries.
We demonstrate how query containment can be decided for federated SPARQL queries under bag and bag-set semantics, opening the door to SPARQL query result caching and reuse, with direct applications to query optimization, incremental querying, and reproducibility of federated SPARQL data access workflows.
\end{abstract}

\maketitle

\section{Introduction}

Minimizing redundant computation is a fundamental concern in query optimization.
Query containment is a classic tool for this purpose: it determines whether the results of one query are always included in those of another, independently of the targeted database.
Solving this problem underpins a wide range of optimization techniques, including caching, view selection, query rewriting, and equivalence validation~\cite{Afrati2019-ch1, DOAN201221, Chaudhuri1993}.
Under set semantics, multiple query fragments have been thoroughly characterized~\cite{Chandra1977, aho1979equivalences, aho1979efficient, Chekuri2000}.
Bag-semantic and bag-set semantic containment have also been studied~\cite{Chaudhuri1993, Afrati2010, khamis2021bag}; however, several query fragments remain undecidable under these semantics and open problems persist.

\ac{KG} has become a widely adopted paradigm for representing and integrating heterogeneous data, enabling applications in scientific data management, enterprise knowledge representation, and the Web of Data.
SPARQL is the standard query language for \ac{KG} built on \ac{RDF}, and federated SPARQL extends this by enabling virtual data integration across distributed \ac{KG} sources, allowing heterogeneous data to be queried jointly without centralization, which is particularly valuable when sources are governed by different policies, maintained by independent parties, or too large to consolidate.
However, federated queries impose significant computational demands on both endpoints and clients, a challenge further compounded by the reliability limitations of public SPARQL endpoints, which frequently experience downtime or performance degradation under load~\cite{buil2013sparql}.
In settings such as scientific reproducibility studies, where queries must be re-executed or iteratively refined, this overhead is especially costly.
Query containment is therefore a natural optimization: when a new query is contained in a previously executed one, results can be derived from cached data without re-querying the federation, reducing both endpoint load and network traffic.

For SPARQL, existing research has addressed set-semantic query containment~\cite{WudageChekol2013}.
Yet the SPARQL query language operates under bag-set semantics~\cite{Perez2006, Angles2016}, in which the underlying database operates under set semantics~\cite{w3ConceptsAbstract} while query operators produce bags\footnote{A bag, or multiset, is similar to set but allows duplicate elements.}~\cite{w3SPARQLQuery}.
Set-semantic containment does not imply bag-set semantic containment, making existing containment results invalid for real-world SPARQL query optimization.

To the best of our knowledge, bag-set semantic SPARQL query containment remains unexplored, as does query containment in the context of federated queries.\footnote{Note that QCan~\cite{salas2018canonicalisation} deals with bag-set semantics, however, it addresses query equivalence rather than query containment, by transforming SPARQL queries into a canonical form and performing syntactic comparison.}

This paper bridges these gaps by formalizing query containment for conjunctive SPARQL federated queries under bag-set semantics.
We first present the \hyperref[sec:related_works]{related works},
then the \hyperref[sec:preliminaries]{preliminaries},
the \hyperref[sec:semantic_federated_q]{semantics of federated queries},
followed by the evaluation of \hyperref[sec:federated_queries_query_containment]{federated queries query containment}
and we \hyperref[sec:conclusion]{conclude} the paper.

\section{Related Work}\label{sec:related_works}

\subsection{Query Containment}\label{sec:related_works_qc}

Query containment is the problem of deciding, for every possible database, whether the result of a query $Q_1$ is a subset of the result of a query $Q_2$, denoted $Q_1 \sqsubseteq Q_2$; in this case $Q_1$ is the \emph{contained query} and $Q_2$ the \emph{containing query}.
Query containment has been used in multiple query optimization tasks, including view selection, caching, and query rewriting~\cite{Afrati2019-ch2, Shirkova2011, Halevy2001}.

Query containment under bag or bag-set semantics is a strictly stronger relationship than its set-semantic counterpart~\cite{Chaudhuri1993}.
The literature does not describe general necessary conditions for containment under bag-set and bag semantics~\cite{ioannidis1995containment, khamis2021bag}; however, sufficient conditions have been identified.
Furthermore, the general case of bag or bag-set semantic query containment for \ac{UCQ} has been shown to be undecidable~\cite{ioannidis1995containment, Jayram2006, Afrati2010}.
However, some fragments of bag-set and bag semantic conjunctive queries have been proven decidable~\cite{khamis2021bag, KOPPARTY20111097, Chaudhuri1993, Afrati2010}.
\Cref{tab_decidable_containment_bag_bag_set} summarizes the bag and bag-set \ac{CQ} containment fragments relevant to this paper and their complexity.\footnote{For the general case, containment mappings give a sufficient but not a necessary condition.}

\begin{table}
  \caption{Complexity of bag and bag-set \ac{CQ} containment for the fragments relevant to this paper.}
  \label{tab_decidable_containment_bag_bag_set}
  \begin{tabular}{lll}
    \toprule
    Case & Semantic & Complexity\\
    \midrule
    General \ac{CQ}~\cite{Chaudhuri1993} & Bag, bag-set & $\Pi^p_2$-hard\\
    \ac{CQ}, $Q_1$ proj.-free~\cite{Afrati2010} & Bag-set & NP-complete\\
    \ac{CQ}, $Q_1,Q_2$ proj.-free~\cite{Afrati2010} & Bag-set & $O(n^2\log n)$\\
    \ac{CQ}, $Q_1,Q_2$ proj.-free~\cite{Afrati2010} & Bag & $O(n^2\log n)$\\
    \bottomrule
  \end{tabular}
\end{table}

Beyond relational databases, query containment has been studied in other graph query languages, including XPath~\cite{schwentick2004xpath} and Cypher~\cite{tang2025proving}.
In the context of SPARQL, set-semantic query containment has received considerable attention~\cite{WudageChekol2013, letelier2013static, chekol2012sparql, chekol2012sparql2, jenaQC, Spasic2023}.
Extensions incorporating RDF schemas, OWL entailment regimes~\cite{WudageChekol2013} and property paths~\cite{Kostylev2025, chekol2016containment} have also been proposed, though these fall outside the scope of this paper.


\subsection{Federated Queries}\label{sec:related_works_fed}

Federated SPARQL queries have been approached from two perspectives: explicit federation through the \texttt{SERVICE} clause, and automatic source selection.
The SPARQL 1.1 Federated Query specification~\cite{w3SPARQLFederated} introduces the \texttt{SERVICE} clause as the W3C standard mechanism for directing subexpressions to specific endpoints.
The SPARQL specification mentions the use of variables as endpoint \acp{IRI} in \texttt{SERVICE} clauses but does not give an official interpretation of their behavior.\footnote{\url{https://www.w3.org/TR/sparql11-federated-query/\#variableService}}
\citeauthor{BuilAranda2013}~\cite{BuilAranda2013} introduced \emph{Service-Safeness}, a syntactic rule outside the SPARQL specification that makes such queries computable by guaranteeing a finite set of federation members.
Because Service-Safeness is not part of the specification and multiple SPARQL query engines use their own interpretation, this paper does not consider variables as endpoint \acp{IRI}.
\citeauthor{Cheng2021}~\cite{Cheng2021} proposed FedQPL, a formalism to express federated queries under automatic source selection.
Automatic source selection algorithms\rt{Let's cite them here! E.g. FedX, ... I'd just include all algorithms that were mentioned and cited in the paper of Elias. We can also include Comunica in this list.} alleviate the practical difficulty of determining which federation member each subquery should be sent to~\cite{hanski2025observations}.
Thus, users can query a federation without knowing how the data is distributed.
Furthermore, \citeauthor{Cheng2021}~\cite{Cheng2021} introduce \ac{ESA}, a canonical source-selection algorithm that requires no prior information about the federation.
In \ac{ESA}, each triple pattern is evaluated against every federation member and the results are combined via bag union, guaranteeing complete results regardless of how data is distributed across the federation.

To the best of our knowledge, no contribution has addressed query containment for federated SPARQL queries; this is the gap our paper addresses.

\section{Preliminaries}\label{sec:preliminaries}

\subsection{SPARQL Language}

\ac{RDF}~\cite{w3ConceptsAbstract} is a data model for representing knowledge on the Web.
Information is expressed as \emph{triples} of the form $(s, p, o)$, where $s$ is the \emph{subject} (the entity being described), $p$ is the \emph{predicate} (the property or relationship), and $o$ is the \emph{object} (the value or related entity).
A set of such triples forms an \ac{RDF} \ac{KG}, which can be thought of as a labeled directed graph: nodes are subjects and objects, and edges are predicates~\cite{w3ConceptsAbstract}.

\acp{IRI} (denoted $\mathcal{I}$) are globally unique identifiers used to name resources and properties.
Literals (denoted $\mathcal{L}$) are concrete data values such as strings, numbers, or dates.
Blank nodes (denoted $\mathcal{B}$) are anonymous resources with no global identifier.
A triple is formally defined as $(s, p, o)$ where $s \in \mathcal{I} \cup \mathcal{B}$, $p \in \mathcal{I}$, and $o \in \mathcal{I} \cup \mathcal{B} \cup \mathcal{L}$.
An \ac{RDF} \ac{KG} is a set of such triples.

SPARQL~\cite{w3SPARQLQuery} is the W3C standard language to query \ac{RDF} \ac{KG}.
SPARQL queries are matched against an \ac{RDF} graph by substituting variables for components of triples.
A \emph{triple pattern} extends a triple by allowing variables (drawn from a set $\mathcal{V}$, disjoint from $\mathcal{I}$, $\mathcal{B}$, and $\mathcal{L}$) in any position:
$\mathit{tp} = (s, p, o)$ where $s \in \mathcal{I} \cup \mathcal{B} \cup \mathcal{V}$, $p \in \mathcal{I} \cup \mathcal{V}$, and $o \in \mathcal{I} \cup \mathcal{B} \cup \mathcal{L} \cup \mathcal{V}$.

A \ac{BGP} is a set of triple patterns evaluated conjunctively, binding variables across all patterns.
A query built from \acp{BGP} is called a \ac{CQ}.
\acp{UCQ} are formed by taking the union of multiple \acp{GP} via a \ac{UGP}, allowing a query to match any of several alternative patterns.

Matching a graph pattern against an \ac{RDF} graph produces \emph{solution mappings}, partial functions from variables to \ac{RDF} terms.
The following definitions make precise two notions on solution mappings used throughout the paper.

\begin{definition}[Multiplicity]\label{def_multiplicity}
  Given a multiset or sequence $\Omega$ of solution mappings and a solution mapping $\mu$, we write $\operatorname{multiplicity}(\mu \mid \Omega)$ to denote the number of times $\mu$ appears in $\Omega$.
\end{definition}

\begin{definition}[Merge~\cite{w3ConceptsAbstract}]\label{def_merge}
  Two solution mappings $\mu_1$ and $\mu_2$ are \emph{compatible} if $\mu_1(v)$ and $\mu_2(v)$ are the same \ac{RDF} term~\cite{w3ConceptsAbstract} for every variable $v \in \operatorname{dom}(\mu_1) \cap \operatorname{dom}(\mu_2)$.
  The \emph{merge} of two compatible solution mappings is their union, written $\operatorname{merge}(\mu_1, \mu_2) = \mu_1 \cup \mu_2$.
\end{definition}

\subsection{Federated Queries}

A SPARQL federated query involves executing different segments of a query across multiple endpoints.
We denote the evaluation of a query over a \ac{KG} $G$ as $\eval{Q}{G}$, and the evaluation over a federation of endpoints $F$ as $\eval{Q}{F}$, where $F$ is defined as a set of federation members, formalized as a mapping $\mathit{fm}\colon \mathit{IRI} \mapsto G$.
Evaluating a query over $F$ involves decomposing $Q$ into subqueries $Q_i$, where each subquery $Q_i$ is evaluated on a knowledge graph $G_j \in \operatorname{codomain}(F)$ member of the federation.
We adopt a database-agnostic view of federation members: since query containment does not depend on the content of any particular knowledge graph, we formalize the notion of federation in terms of endpoint \acp{IRI} rather than data.

We distinguish three types of federated queries: \ac{UDF}, \ac{ESA}~\cite{Cheng2021}, and \ac{ASSF}.
In \ac{UDF}, the query specifies \texttt{SERVICE} clauses~\cite{w3SPARQLFederated} that explicitly name the target endpoint for each subquery; a variable may be used as the endpoint \ac{IRI}, in which case the query must satisfy \emph{Service-Safeness}~\cite{BuilAranda2013} to guarantee a finite federation.
In \ac{ESA}, a finite federation $F$ is provided before execution and each triple pattern is evaluated over every $G \in \operatorname{codomain}(F)$, with results accumulated via bag union~\cite{w3SPARQLFederated}.
In \ac{ASSF}, a finite $F$ is similarly given, but the query engine selects which federation member receives each triple pattern rather than broadcasting to all.

\begin{definition}[FedQPL operators~\cite{Cheng2021}]\label{def_fedqpl}
  We express federated query plans with three operators of the FedQPL formalism.
  The \emph{request} $\freq_{\mathit{fm}}^{P}$ evaluates a graph pattern $P$ against a single federation member $\mathit{fm}$.
  The \emph{multiway union} $\fmu$ and \emph{multiway join} $\fmj$ combine the results of their arguments by bag union and by join, respectively.
\end{definition}

\section{Semantics of Federated SPARQL Queries}\label{sec:semantic_federated_q}

Before investigating query containment, we must establish under which semantics the different queries we study are evaluated.
This choice determines whether and how duplicate results can arise, which directly impacts how query containment is resolved.

\subsection{Single Endpoint Execution}\label{sec:single_endpoint}
SPARQL queries on single sources are evaluated under bag-set semantics.
The underlying \ac{KG} operates under set semantics~\cite{w3ConceptsAbstract}, whereas SPARQL operators produce bags.
The choice between bag and set semantics has a significant impact on the complexity and decidability of query containment~\cite{Chaudhuri1993}.
For relational database settings, \citeauthor{Chaudhuri1993}~\cite{Chaudhuri1993, Afrati2010} established that for \acp{CQ} only projection can increase the multiplicity of a result.
The same property holds for SPARQL \acp{CQ}, where the only body operator, \texttt{JOIN}, cannot raise the multiplicity, as in the relational case~\cite{Chaudhuri1993, Afrati2010}; projection at the head therefore remains the only operator that can do so.
In the context of \ac{UCQ}, the multiplicity can additionally be increased by \texttt{UNION} and by the combination of \texttt{UNION} and \texttt{JOIN}, as defined in \Cref{def_union_multiplicity,def_join_multiplicity}.
\ac{UCQ} is not the subject of our study; however, in the next sections we introduce \ac{UCFQ}, a subclass of \ac{UCQ}, together with the semantics of federated queries.



\subsection{Federated Queries}
A \ac{KG} is stored under set semantics, and querying a single endpoint operates under bag-set semantics.
A federated query, however, combines solution mappings drawn from several \acp{KG} at once, so it is not obvious whether spanning multiple members can introduce duplicate results.
We therefore examine the semantics of federated query evaluation.
The SPARQL federated query specification~\cite{w3SPARQLFederated} defines the execution of federated queries with \texttt{SERVICE} clauses as a join between the bag of solution mappings from the preceding execution and the bag of solution mappings returned by the SPARQL service.
The \ac{GP} inside a \texttt{SERVICE} clause has no projection, so each federation member returns a set of solution mappings.
A \texttt{SERVICE} clause is nested in the query through a \texttt{JOIN}, which cannot increase the multiplicity.
Hence, even when some federation members hold duplicate information, the federated query produces no duplicate results, and \ac{CQ} evaluation with \texttt{SERVICE} clauses remains under bag-set semantics.

\subsection{\acl{UCFQ}}\label{sec:ucfq}

We define a particular class of federated queries, useful for capturing duplicate information held by different federation members.
Each triple pattern of the query body is evaluated across several members and the partial results are joined.

\begin{definition}[\acl{UCFQ}]\label{def_ucfq}
A \ac{UCFQ} is a multiway join $\fmj(U_1, \dots, U_n)$ where each operand $U_i = \fmu(\freq_{I_1}^{\mathit{tp}_i}, \dots, \freq_{I_{k_i}}^{\mathit{tp}_i})$ is a multiway union evaluating the same triple pattern $\mathit{tp}_i$, up to a renaming of the variables local to $U_i$, at a subset of federation members $I_1, \dots, I_{k_i}$, optionally followed by a projection.
\end{definition}

A variable is \emph{local} to $U_i$ when the head discards it and no other operand of the join binds it.
Renaming the local variables of a branch, distinct variables to distinct variables, maps its solution mappings one to one and preserves their multiplicity, and no other operator of the query binds them, so the result bag is unchanged.
Branches naming them differently therefore evaluate the same triple pattern, which matters as a query engine binding a variable at a single branch may name it freely.

A \ac{CQ} is the special case where every union ranges over a single member, and every \ac{UCFQ} is a \ac{UCQ}, so the \acp{UCFQ} lie strictly between the \acp{CQ} and the \acp{UCQ}.
In particular, each triple pattern of a \ac{CQ} is the single-member union $\fmu(\freq_{\mathit{fm}}^{\mathit{tp}})$ over the member $\mathit{fm}$ at which it is evaluated.

\acp{UCFQ} arise under automatic source selection, in particular \ac{ESA} (\Cref{sec:related_works_fed}).
Consider a federation of two members, a job registry and a social network, that redundantly store the job and birth date of a person.
The projection-free \ac{CQ} retrieving each person's job and birth date (\Cref{fig:ucfq-example}, top) would return a set on a single endpoint, but under \ac{ESA} it is evaluated as the \ac{UCFQ} (\Cref{fig:ucfq-example}, bottom), unioning each triple pattern across both members before joining the results.
Since both members hold the same solution for a person, the join returns that person with multiplicity four, even though every member stores a set --- a duplication that does not arise with explicit \texttt{SERVICE} clauses in a \ac{CQ}.%
\footnote{An executable example of the semantics of federated queries, evaluated with the Comunica~\cite{taelman_iswc_resources_comunica_2018} query engine, is available at \url{https://github.com/constraintAutomaton/demo-federated-query-semantic}.}

\begin{figure}
  \textbf{The \ac{CQ} as written}
  \lstinputlisting[style=sparql]{code/ucfq_cq.rq}
  \textbf{The executed query under \ac{ESA}}
  \lstinputlisting[style=sparql]{code/ucfq_example.rq}
  \caption{A \ac{CQ} evaluated as a \ac{UCFQ} under \ac{ESA}.}
  \label{fig:ucfq-example}
\end{figure}

\ac{ESA} is the particular case of \ac{UCFQ} in which every union branch ranges over the entire federation $\{\mathit{fm}_1, \dots, \mathit{fm}_m\}$, formalized as follows:
\begin{equation}
  \fmj\bigl(\fmu(\freq_{\mathit{fm}_1}^{\mathit{tp}_1}, \dots, \freq_{\mathit{fm}_m}^{\mathit{tp}_1}), \dots, \fmu(\freq_{\mathit{fm}_1}^{\mathit{tp}_n}, \dots, \freq_{\mathit{fm}_m}^{\mathit{tp}_n})\bigr)
  \label{eq_esa}
\end{equation}

Most federated engines produce plans of this same join-over-union shape---for instance FedX~\cite{schwarte2011fedx}, SemaGrow~\cite{charalambidis2015semagrow}, and CostFed~\cite{saleem2018costfed}---which can likewise be rewritten as \acp{UCFQ}.
Automatic source selection can, however, produce \ac{CQ} plans of arbitrary shape --- for instance, union-over-join plans~\cite{aimonier2024fedup} --- and \ac{UCFQ} therefore does not capture every \ac{CQ} under automatic source selection.

Containment of \acp{UCQ} under bag-set semantics is undecidable (\Cref{sec:related_works_qc}), but this does not settle the status of the \acp{UCFQ}, which form a strict subclass.
The next section introduces the \ac{BFR}, a transformation showing that evaluating a \ac{UCFQ} is equivalent to evaluating a \ac{CQ} over one \ac{BKG} per union branch under bag semantics.



\subsection{\acl{BFR}}\label{sec_bag_federation_reduction}

We propose the \ac{BFR}, a transformation that rewrites a \ac{UCFQ} as a \ac{CQ} evaluated over one \ac{BKG} per union branch (\Cref{def_bkg}).
Intuitively, a virtual operator aggregates the members of each branch's sub-federation into a \ac{BKG} and evaluates the branch's triple pattern over it.

\subfile{proofs/bag_federation_reduction}

\subsection{A Bag-Semantic Reading of FedQPL}\label{sec_fedqpl}

\Cref{sec:single_endpoint,sec:ucfq} established that SPARQL evaluates under bag-set semantics: the underlying \acp{KG} are sets, while the operators combining their results return bags, so that unioning a triple pattern across federation members duplicates solution mappings (\Cref{def_union_multiplicity}).
We expressed the queries this produces in the notation of FedQPL (\Cref{def_ucfq}), which carries no such semantics: \citeauthor{Cheng2021} evaluate every operator over sets of solution mappings, and $\fmu$ denotes set union~\cite[Def.~6]{Cheng2021}.
We read $\fmu$ and $\fmj$ as their bag counterparts (\Cref{def_fedqpl}), which aligns them with the SPARQL operators they lift and places FedQPL under the bag-set semantics of \Cref{sec:single_endpoint}.
Other readings are possible, and we return to them below.
\citeauthor{Cheng2021} give a bag semantics as future work~\cite{Cheng2021}, and to the best of our knowledge no such extension has been published; we set out here what one has to settle.

The federation semantics of FedQPL evaluates a \ac{GP} over the set union $G_{\mathit{union}}$ of the \acp{KG} of the federation members~\cite[Def.~3]{Cheng2021}, following the stated intention that ``the result of a SPARQL query over such a federation should be the same as if the query was executed over the union of all the RDF data available in all the federation members''.
Correctness of a source assignment is defined against that graph~\cite[Def.~7]{Cheng2021}.
Under set semantics the exhaustive source assignment is correct~\cite[Prop.~1]{Cheng2021}; under bag-set semantics it is not.

\begin{proposition}\label{prop_fedqpl_esa_incorrect}
  Read under bag-set semantics, the exhaustive source assignment~\cite[Def.~9]{Cheng2021} is not correct in the sense of~\cite[Def.~7]{Cheng2021}.
\end{proposition}

\begin{proof}
  $G_{\mathit{union}}$ is a set, so a \ac{BGP} returns each of its solution mappings once over it (\Cref{sec:single_endpoint}).
  The \ac{ESA} expression of \Cref{eq_esa}, on the other hand, returns $\mu$ with multiplicity $\prod_{i} |\{\mathit{fm} \in F : \mu(\mathit{tp}_i) \in G_{\mathit{fm}}\}|$ by \Cref{def_union_multiplicity,def_join_multiplicity}.
  The federation of \Cref{fig:ucfq-example} makes this product four, whereas $\eval{B}{G_{\mathit{union}}}$ returns that solution mapping once.
\end{proof}

There are two ways to preserve~\cite[Def.~3]{Cheng2021} as stated.
The first assumes that no triple is duplicated across the federation members, so that the set union of their \acp{KG} coincides with the bag union.
Such an assumption is hard to verify in the open-Web setting and cannot always be enforced.
The second introduces a federation-level union operator that deduplicates the contributions of distinct members.
That operator is not a set union, but a different and more complex type of union.
A request may itself return a bag, since the request language of a SPARQL endpoint is the full graph pattern language, and those multiplicities must survive.
The operator must therefore deduplicate across members while leaving the multiplicities of each member's own response untouched, which requires the algebra to track the member every solution mapping comes from --- an operator SPARQL does not have, and whose addition changes the algebra rather than extending it.
We adopt neither.

We propose instead to use a \ac{BKG} (\Cref{def_bkg}) in place of the set union.
The federation is then a bag rather than a set, so the reading moves from bag-set to bag semantics, which is the transition the \ac{BFR} performs (\Cref{sec_bag_federation_reduction}).
\Cref{prop_duplicate_fed_mem}, taken with $F_i = F$ for every branch, restores the correctness of the exhaustive source assignment, and the intention of \citeauthor{Cheng2021} survives with the set union of the \acp{KG} replaced by their bag union.
This reading keeps the \texttt{UNION} of SPARQL and adds no operator to the algebra.
It also leaves the duplication of a triple across federation members observable in the result bag, so that a query can determine how many members hold a given triple --- information that serves domain purposes as well as the checking of integrity constraints over a federation.
We adopt it for the remainder of the paper.

\section{Conditions for Federated Query Containment}\label{sec:federated_queries_query_containment}

Federated query containment may initially appear fundamentally different from single-source query containment, since containment is data-source independent while a federated query names particular databases.
However, this difference is not fundamental.
We propose that federated query containment must consider the \acp{IRI} \rt{Important!!!: Don't use IRI in this paper, but use URL, as it more widely understood.} of the federation members, but not their content.\rt{Might be good to add a footnote or so to mention that we assume immutable data in federation members. Because if you'd do caching for example, and the endpoints change their data, you need more than just the URL, which is out of scope for this work here.}
Referencing these \acp{IRI} could benefit the resolution of query containment, for instance by guaranteeing the absence of overlap between members, since bag semantics is generally more complex to decide than bag-set semantics.
We leave such optimizations out of scope.

Without loss of generality, we work only with flattened \texttt{SERVICE} clauses, where every \texttt{SERVICE} clause appears at the \ac{BGP} level.
The proposition below shows that a nested \texttt{SERVICE} clause is semantically equivalent to such a flattened form, so nesting reduces to this case.\footnote{Nesting is also discouraged in practice\rt{There are no official statements or papers AFAIK that discourage this, so please rephrase this, as it's not accurate.}, as delegating the remote execution of a \texttt{SERVICE} clause to a third-party engine can raise load-management concerns and leaves the query issuer unable to guarantee that the endpoint supports federated query execution.}

\subfile{proofs/nested_service_flattening}


\subsection{Query Containment Condition}~\label{sec:federated_queries_query_containment_condition}

\rt{The purpose of this section is unclear; it's not explained. Why are the following functions introduced?}

We define a function $\mathit{QS}\colon (\mathcal{Q}, \mathcal{F}) \to \mathcal{Q}$ that returns the query associated with the \ac{GP} of a \texttt{SERVICE} clause.
Here $\mathcal{F}$ is the set of all federation members and $\mathcal{Q}$ is the set of all queries.
We also define $\mathit{QO}\colon \mathcal{Q} \to \mathcal{Q}$, a function that returns the segments of a query evaluated against the local \ac{KG}.

\subfile{proofs/federated_query_condition}

\section{Conclusion}\label{sec:conclusion}

This paper formalizes query containment for conjunctive SPARQL federated queries under bag-set semantics.
We distinguish two source-selection mechanisms, \ac{UDF} and \ac{ASSF}, and an orthogonal axis between \ac{SSA} and \ac{DSA} that captures when source assignment is determined.
We introduce the \ac{BFR}, a transformation showing that \ac{ASSF} evaluation is equivalent to evaluating the original \ac{CQ} against a \ac{BKG} under bag semantics, which lets us inherit existing bag-semantic \ac{CQ} containment results into the federated setting.

These results open the door to result caching and reuse across federated SPARQL workflows.
This has direct applications to query optimization, incremental querying, and the reproducibility of federated SPARQL data access workflows.

Future work could consist of investigating containment between queries evaluated under bag and bag-set semantics in the context of the \ac{BFR}, analyzing the complexity of the proposed containment condition, implementing it in a query containment solver, and extending the framework to additional SPARQL fragments.


\bibliographystyle{ACM-Reference-Format}
\bibliography{references}

\end{document}
